\documentclass[11pt,a4paper]{article}
\usepackage[margin=19mm,headheight=15pt]{geometry}
\usepackage{fontspec}
\setmainfont{Noto Serif}
\usepackage{amsmath,amssymb,mathtools}
\usepackage{xcolor}
\usepackage{fancyhdr}
\usepackage{enumitem}
\usepackage{microtype}
\setlength{\parindent}{0pt}
\setlength{\parskip}{6pt}
\definecolor{navy}{HTML}{183153}
\definecolor{lightgray}{HTML}{F3F5F7}
\pagestyle{fancy}\fancyhf{}
\lhead{\small ME 5309 -- Mid-Semester Examination}
\rhead{\small ME 5309}
\cfoot{\thepage}
\begin{document}
\begin{center}
{\Large\bfseries ME 5309 -- Mid-Semester Examination}\\[3pt]
{\large Answer Booklet}\\[8pt]
\end{center}
\textbf{Instructions:} Attempt any three questions. Each question carries 10 marks. Show working where appropriate.
\section*{Question 1 \hfill [10 marks]}
\subsection*{Q1(a) \hfill [4 marks]}

The pixel pitch is 0.005 mm, so f\_x = 16/0.005 = 3200 pixels.
Sensor width = 6.4 mm.
Belt field of view = 400x6.4/16 = 160 mm; spatial resolusion = 160/1280 = 0.125 mm/pixel.

\subsection*{Q1(b) \hfill [3 marks]}

ADC step = 5/4096 = 1.2207 mV.
I divide by 0.1 mV/N, giving ΔF = 12.207 N.

\subsection*{Q1(c) \hfill [3 marks]}

For zero-order hold, F = e\textasciicircum{}(AT) and G = integral\_0\textasciicircum{}T e\textasciicircum{}(Aτ)B dτ.
Here:\[\begin{aligned}A = -2, B = 1, T = 0.1 s\end{aligned}\]
Thus:\[\begin{aligned}F = e^(-0.2) = 0.818731\\ G = (1-e^(-0.2))/2 = 0.0906346\end{aligned}\]
Therefore:\[\begin{aligned}x_{k+1} = 0.818731 x_k + 0.0906346 u_k\end{aligned}\]\par\smallskip With $AT=(-2)(0.1)=-0.2$, $F=e^{-0.2}=0.8187308$. For $G$, $(1-F)/2=(1-0.8187308)/2=0.1812692/2=0.0906346$. Thus the input contribution is positive.

\section*{Question 3 \hfill [10 marks]}
\subsection*{Q3(a) \hfill [2 marks]}

First, numerical differentiation amplifies high-frequency encoder noise and quantisation because a small position error is divided by the short sample interval.
Second, differencing is sensitive to sampling/timing and can produce a noisy, jittery velocity signal; filtering reduces noise but adds lag.\par\smallskip In a finite difference, the change in position is divided by the sample interval. Encoder quantization error appears in the difference of two measurements, so it can be magnified by the division by $T_s$. The result may jump around even if the actual motion is smooth.

\subsection*{Q3(b) \hfill [2 marks]}

\[\begin{aligned}y[k]=0.8818y[k-1]+0.0591x_k+0.0591x[k-1], DC gain 1\end{aligned}\]

\subsection*{Q3(c) \hfill [4 marks]}

A-LC=[[-ℓ₁,1], [-ℓ₂,-2]].
I obtain p(s)=s²+(ℓ₁+2)s+(ℓ₁+ℓ₂).
Matching (s+10)² gives ℓ₁=18 and ℓ₂=82.

\subsection*{Q3(d) \hfill [2 marks]}

The observer is simpler because it only smooths the measured velocity.
The filter is safer when the model is poor.

\section*{Question 4 \hfill [10 marks]}
\subsection*{Q4(a) \hfill [2 marks]}

The characteristic polynomial is s²+4s+8, so I report poles -4 and -8.
Unity DC gain requires Nff=8.

\subsection*{Q4(b) \hfill [2 marks]}

With actual b=0.75 and unchanged Nff=8, k₁=2, the steady-state relation is θss = bNff/(c+bk₁) = 0.75x8/(6+0.75x2)=6/7.5=0.8.
For r=1, percentage error = (1-0.8)x100\%=20\%.\par\smallskip With the actual gain, the steady-state output is $y_{ss}=b_{actual}N_{ff}r/ (c+b_{actual}k_1)$ for the stated setup. Substituting the given numbers gives $y_{ss}=0.8$ for a unit reference. Therefore $e_{ss}=1-0.8=0.2$, i.e. a 20\% error.

\subsection*{Q4(c) \hfill [4 marks]}

I obtain s³+(a+bk₂)s²+(c+bk₁)s-bK\_I.
Desired polynomial is s³+8s²+24s+32, so k₂=6, k₁=18, K\_I=-32.

\subsection*{Q4(d) \hfill [2 marks]}

At steady state, a stable closed loop with integral action drives ξ̇=r-θ to zero; therefore r-θ=0 and the tracking error is zero, assuming no saturation or other implementation limits prevent convergence.
A practical failure mode is integrator windup during actuator saturation; the standard remedy is anti-windup, such as back-calculation or conditional integration.\par\smallskip With integral action, a constant tracking error keeps accumulating in the integral state, so the controller can adjust its output until the error goes to zero, assuming the closed loop is stable and the actuator is not constrained. Under saturation the integral state can keep growing; anti-windup limits or back-calculates that state.

\end{document}